% !TeX spellcheck = ru_RU
% !TEX root = vkr.tex

% добавить полноценные описания двух процессоров и лекции

\section{Описание решения}

% !TeX spellcheck = ru_RU
% !TEX root = vkr.tex

\subsection{Технологии и инструменты разработки}

Для проектирования виртуальных вычислительных устройств выбран язык описания аппаратуры \textbf{Verilog HDL}. Verilog широко применяется для моделирования цифровых систем различной сложности благодаря синтаксису, близкому к языку C, что упрощает освоение. Язык поддерживает структурное, поведенческое и потоковое описание, обеспечивая гибкость при создании логических схем и архитектурных решений. Симуляция Verilog-моделей выполняется с помощью \textbf{Icarus Verilog} -- свободного инструмента с поддержкой компиляции, интерактивной отладки и высокой производительностью. 

Для системного программирования используется язык \textbf{C}. Компиляция программ под целевую архитектуру RISC-V выполняется кросс-компилятором \textbf{riscv64-unknown-elf-gcc}. Вспомогательные скрипты автоматизации написаны на \textbf{Python 3}. Компиляция ПО осуществляется компиляторами \textbf{GCC} и \textbf{TCC} с поддержкой архитектуры RISC-V. В более сложных проектах для изоляции окружения и гарантии воспроизводимости сборки применяется Docker, который включает в себя RISC-V тулчейн, Verilator, SystemC и другие зависимости.

\subsection{Моделирование и эволюция систем}

Разработка виртуальных устройств выполняется поэтапно, что соответствует дидактическому принципу «от простого к сложному» и позволяет последовательно осваивать ключевые концепции computer science:

\begin{enumerate}
  \item \textbf{Минималистичная модель} (NAND\_CPU) -- ограниченный набор команд для изучения базовых принципов вычислений.
  \item \textbf{Расширенные процессоры} (riscv32\_cpu\_project) -- полный набор базовых инструкций, таких как RISC-V RV32I, для выполнения типовых вычислительных задач. Примером служит реализация однотактового ядра.
  \item \textbf{Сложные системы} (verilog-c-os) -- платформы, поддерживающие выполнение программ на языке C и операционных систем. Проект verilog-c-os, включающий версии для запуска «голого» C-кода и учебной ОС xv6. Такие системы демонстрируют взаимодействие аппаратного и программного уровней, включая управление процессами и памятью.
\end{enumerate}

Поэтапный подход обеспечивает управляемость разработки и постепенное наращивание функциональности.

\subsection{Структура и архитектура виртуальных устройств}

\subsubsection{NAND\_CPU}
Базовая архитектура состоит из ключевых компонентов:
\begin{itemize}
  \item \textbf{Регистровый файл} -- асинхронное чтение, синхронная запись;
  \item \textbf{Память} -- унифицированное хранение инструкций и данных;
  \item \textbf{ALU} -- арифметико-логические операции под управлением кодов команд;
  \item \textbf{Блок управления} -- декодирование и формирование сигналов синхронизации.
\end{itemize}

Инструкции исполняются пошагово с поддержкой условных переходов.

\subsubsection{RISC-V RV32I CPU}
Данная реализация представляет собой \textbf{синтезируемое, однотактовое (single-cycle) ядро}, разработанное как \textbf{образовательный эталон}. Оно полностью поддерживает базовый набор целочисленных инструкций RISC-V (RV32I). Процессор выполняет одну инструкцию за такт; конвейер, кэши и механизмы предсказания переходов отсутствуют, что делает тракт данных полностью прозрачным для анализа и соответствует критериям учебного компьютера, перечисленным в обзоре.

\textbf{Набор инструкций и архитектура:}
Поддерживаются все основные форматы инструкций RV32I, что позволяет выполнять любые целочисленные вычисления и управляющие конструкции:

\begin{center}
\begin{tabular}{|l|>{\ttfamily}p{5cm}|p{8cm}|}
\hline
\textbf{Формат} & \multicolumn{1}{c|}{\textbf{Примеры команд}} & \textbf{Назначение} \\
\hline
\textbf{R-type} & ADD, SUB, AND, OR, XOR, SLT, SLL & Операции между регистрами \\
\hline
\textbf{I-type} & ADDI, LW, JALR & Операции с константами, загрузка, косвенный переход \\
\hline
\textbf{S-type} & SW & Сохранение в память \\
\hline
\textbf{B-type} & BEQ, BNE, BLT, BGE, BLTU, BGEU & Условные переходы (ветвления) \\
\hline
\textbf{J-type} & JAL & Безусловный переход и сохранение адреса возврата \\
\hline
\textbf{U-type} & LUI, AUIPC & Загрузка старших 20 бит константы \\
\hline
\end{tabular}
\end{center}

Нераспознанные коды операций активируют сигнал \texttt{trap}, перенаправляя выполнение по вектору прерывания (\texttt{0x00000010}), что является базой для реализации исключений.

\textbf{Наличие и тип компонентов:}
Архитектура включает специализированные модули Verilog, каждый из которых отвечает за строго определенную функцию:

\begin{itemize}
    \item \texttt{ALU.v} – арифметико-логическое устройство (сложение, вычитание, логические операции, сдвиги);
    \item \texttt{ALUControl.v} – декодер операций ALU по сигналам \texttt{alu\_op}, \texttt{funct3} и \texttt{funct7};
    \item \texttt{BranchComparator.v} – модуль проверки всех шести условий ветвлений RV32I (равенство, меньше, меньше или равно);
    \item \texttt{ControlUnit.v} – центральный декодер, формирующий управляющие сигналы на основе кода операции (opcode);
    \item \texttt{DataMemory.v} – синхронная память данных (RAM на 256 слов) с записью по тактовому сигналу;
    \item \texttt{ImmediateGenerator.v} – формирование знаково-расширенных непосредственных значений для форматов I/S/B/U/J;
    \item \texttt{InstructionMemory.v} – комбинационная память инструкций (ROM на 256 слов), инициализируемая из hex-файла;
    \item \texttt{PCUpdate.v} / \texttt{ProgramCounter.v} – логика обновления счетчика команд с четким приоритетом: Trap → Jump → Branch → \texttt{PC+4};
    \item \texttt{RegisterFile.v} – файл из 32 регистров (асинхронное чтение, синхронная запись, регистр \texttt{x0} аппаратно привязан к нулю);
    \item \texttt{RV32ICPU.v} – топовый модуль, объединяющий все компоненты и реализующий полный тракт данных за один такт.
\end{itemize}

\textbf{Способ реализации и интерфейс:}
Проект реализован на Verilog HDL и ориентирован на симуляцию в Icarus Verilog. Для взаимодействия предусмотрена пакетная обработка: ассемблерная программа (\texttt{prog.txt}) транслируется Python-скриптом \texttt{encode.py} в hex-файл, который затем загружается в память модели. Для отладки в модуль \texttt{RV32ICPU} встроены блоки с \texttt{\$display}, выводящие в консоль детальную информацию о каждом такте (PC, инструкция, значения регистров, обращения к памяти). Дополнительно генерируется VCD-файл для просмотра временных диаграмм в GTKWave, что обеспечивает высокую интерактивность и наглядность, соответствующую критериям учебного компьютера.

\subsubsection{Платформа для запуска C и ОС}
Этот проект расширяет возможности базового ядра и предоставляет полноценную симуляционную среду для выполнения программ на языке C и учебной операционной системы. Он демонстрирует, как абстракции ОС (управление процессами, памятью) реализуются поверх аппаратной архитектуры.

\textbf{Структура проекта:}
Репозиторий содержит две самостоятельные версии с идентичной внутренней структурой, что упрощает навигацию и изучение:
\begin{itemize}
    \item \textbf{Версия \texttt{run\_c\_version}} – минимальное окружение для запуска «голых» C-программ без операционной системы. Включает загрузочный код (\texttt{startup.s}), скрипт компоновщика (\texttt{link\_script.ld}) и небольшую библиотеку поддержки (\texttt{utils.c}), предоставляющую базовые функции ввода-вывода.
    \item \textbf{Версия \texttt{xv6\_version}} – реализует упрощённое ядро, подобное xv6, для 32-битного RISC-V. Ядро поддерживает ключевые функции ОС, описанные в обзоре: управление процессами (таблица процессов) и базовое выделение памяти.
\end{itemize}

Обе версии имеют одинаковую организацию каталогов:
\begin{itemize}
    \item \texttt{core/riscv/} – Verilog RTL ядра процессора (аналогичное описанному выше, но с поддержкой более сложной периферии).
    \item \texttt{top\_tcm\_axi/} – тестовое окружение верхнего уровня и модель памяти (TCM – Tightly Coupled Memory) с интерфейсом AXI.
    \item \texttt{isa\_sim/} – быстрая модель процессора на C++ (ISA-симулятор) для отладки ПО без детального моделирования тактов.
    \item \texttt{docker\_image/} – Docker-образ со всеми необходимыми инструментами, обеспечивающий воспроизводимость окружения.
    \item \texttt{build/} – директория для сгенерированных артефактов.
\end{itemize}

\textbf{Режимы симуляции и взаимодействие:}
Для каждой версии доступны два пути симуляции, что позволяет выбрать баланс между скоростью и детализацией:
\begin{itemize}
    \item \textbf{ISA-симулятор} (C++ модель) – высокая скорость, подходит для функциональной отладки алгоритмов ядра и прикладных программ.
    \item \textbf{RTL-симулятор} (Verilator + Verilog) – цикл-точная симуляция, генерация VCD-файлов для детального анализа аппаратных сигналов.
\end{itemize}

Результаты симуляции сохраняются в папке \texttt{build/}, предоставляя исчерпывающую информацию для анализа:
\begin{itemize}
    \item \texttt{*.elf}, \texttt{*.disasm} – исполняемый файл и его дизассемблерный листинг.
    \item \texttt{\_isa.log} – лог от быстрой ISA-модели.
    \item \texttt{\_verilog.log} – лог UART от RTL-симуляции (эмуляция последовательного порта).
    \item \texttt{\_verilator.vcd} – файл波形 с детальными сигналами процессора для просмотра в GTKWave.
\end{itemize}

\subsection{Реализация и интеграция программного обеспечения}

Процесс интеграции программного обеспечения с аппаратными моделями унифицирован для обоих проектов и включает следующие этапы, соответствующие критериям учебных компьютеров:

\begin{enumerate}
    \item \textbf{Компиляция} – исходные тексты на C и ассемблере компилируются с помощью \texttt{riscv64-unknown-elf-gcc} в ELF-файлы.
    \item \textbf{Формирование образов памяти} – из ELF-файлов извлекаются машинные коды и данные, которые преобразуются в hex-формат с помощью Python-скриптов (например, \texttt{encode.py} из первого проекта).
    \item \textbf{Загрузка в симулятор} – Verilog-тестбенчи загружают hex-файлы в модули \texttt{InstructionMemory} и \texttt{DataMemory} через системную функцию \texttt{\$readmemh}.
    \item \textbf{Инициализация} – для версий с ОС разработаны загрузчики, выполняющие начальную настройку стека, инициализацию BSS-секции и переход к точке входа ядра (функция \texttt{main}).
\end{enumerate}

Такой модульный подход обеспечивает гибкость, быструю интеграцию компонентов и удобство для учебных целей. Использование Docker гарантирует, что любой студент может воспроизвести окружение на своей машине без сложной настройки инструментов.

\subsection{Лекция «От транзистора до операционной системы»}

В состав образовательного комплекса входит лекция, которая формирует у студентов целостное представление об устройстве компьютера — от физической основы вычислений до работы операционной системы. Лекция построена как последовательное восхождение по уровням абстракции и служит теоретической базой для выполнения практических работ с проектами \texttt{riscv32\_cpu\_project} и \texttt{verilog-c-os}.

\subsubsection{Содержание лекции}

Лекция начинается с объяснения роли транзистора как базового элемента цифровой логики. На простых примерах демонстрируется, как из транзисторов строятся логические вентили (НЕ, И, ИЛИ), и показывается, что вентиль И-НЕ является функционально полным.

Далее рассматривается построение более сложных компонентов:
\begin{itemize}
    \item \textbf{Комбинационные схемы:} полный сумматор, декодер, компаратор.
    \item \textbf{Последовательностные схемы:} регистры, счётчики, принципы организации статической памяти (SRAM).
    \item \textbf{Реализация памяти:} объясняется устройство ячейки флеш-памяти на транзисторе с плавающим затвором — запись, стирание, считывание, понятие многоуровневых ячеек (MLC, TLC, QLC).
\end{itemize}

На следующем этапе лекция переходит к синтезу простейшего процессора: рассматриваются составные части (счётчик команд, память инструкций, декодер, регистровый файл, АЛУ), система команд и цикл выполнения (fetch–decode–execute–writeback).

Затем обсуждаются базовые принципы работы устройств ввода-вывода (клавиатура, экран) и роль компилятора в трансляции пользовательских программ в машинный код.

Центральная часть лекции посвящена операционной системе на примере учебного ядра xv6 для RISC-V. Разбираются ключевые механизмы:
\begin{itemize}
    \item управление процессами (состояния, переключение контекста, планировщик round-robin);
    \item виртуальная память (трансляция адресов через таблицы страниц, обработка page fault);
    \item файловая система (структура диска, иноды, буферный кэш);
    \item системные вызовы как интерфейс между пользовательским кодом и ядром.
\end{itemize}

Лекция завершается сопоставлением учебной модели xv6 с реальными встраиваемыми Linux-системами, обсуждаются компоненты таких систем (загрузчик, ядро, корневая файловая система, тулчейн) и инструменты сборки (Buildroot, Yocto). Важным элементом является перечисление того, что отсутствует в xv6 (пользователи и права, сеть, драйверы, графический интерфейс и др.), что помогает студентам осознать упрощённый характер учебной ОС и понять, какие задачи решаются в промышленных разработках.